\Askhsh{24703}{2}
\begin{ylh}{1}
    \item 1.2 Συναρτήσεις
    \item 1.3 Μονότονες συναρτήσεις - Αντίστροφη συνάρτηση.
\end{ylh}
\begin{wrapfigure}[15]{r}{7.5cm} % Adjust height as needed for text wrapping
    \centering
    \begin{tikzpicture}
        \begin{axis}[
            axis lines = middle,
            xlabel = {}, % No explicit labels x, y
            ylabel = {},
            xmin = -0.5, xmax = 1.2,
            ymin = -0.5, ymax = 1.2,
            xtick = {1},
            ytick = {1},
            grid = none,
            samples = 200, % Higher samples for smoother curves
            width = 7cm, height = 6.5cm, % Adjust size to fit wrapfigure and aspect ratio
            tick label style = {font=\scriptsize},
        ]
            % Plot f(x) = sqrt(1-x)
            \addplot[very thick, maincolor, domain=-0.5:1] {sqrt(1-x)};
            \node[maincolor, anchor=west] at (axis cs:-0.2, 1.1) {$f(x) = \sqrt{1-x}$};

            % Plot f^{-1}(x) = 1-x^2 (the portion shown in the figure)
            \addplot[very thick, secondarycolor, domain=0:1.1] {1-x^2};
            \node[secondarycolor, anchor=north east] at (axis cs:0.7, 0.6) {$C_{f^{-1}}$};

            % Plot y=x
            \addplot[dashed, gray, domain=-0.4:1.1] {x};
            \node[gray, anchor=north west] at (axis cs:1.0, 1.1) {$y=x$};

            % Axis labels O, 1, 1
            \node[below left, inner sep=1pt] at (axis cs:0,0) {O};
            \node[below, inner sep=1pt] at (axis cs:1,0) {1};
            \node[left, inner sep=1pt] at (axis cs:0,1) {1};

            % Points for (0,1), (1,0) (blue dots in original) and intersection point (black dot)
            \addplot[only marks, mark=*, mark size=1.5pt, maincolor] coordinates {(0,1) (1,0)};
            \addplot[only marks, mark=*, mark size=1.5pt, black] coordinates {(0.618,0.618)}; % Intersection of f(x) = x, i.e., sqrt(1-x)=x -> x = (sqrt(5)-1)/2
        \end{axis}
    \end{tikzpicture}
\end{wrapfigure}
Θεωρούμε τη συνάρτηση $f$ με $f(x) = \sqrt{1-x}$ και $x \in (-\infty, 1]$.
\begin{alist}
\item Να αποδείξετε ότι υπάρχει η συνάρτηση $f^{-1}$. \monades{8}
\item Να βρείτε τη συνάρτηση $f^{-1}$. \monades{10}
\item Στο παρακάτω σχήμα δίνεται η γραφική παράσταση της συνάρτησης $f$ και ένα τμήμα της γραφικής παράστασης της $f^{-1}$. Να μεταφέρετε στο φύλλο απαντήσεων το παραπάνω σχήμα και το οποίο να συμπληρώσετε με την υπόλοιπη γραφική παράσταση της συνάρτησης $f^{-1}$. \monades{7}
\end{alist}